\documentclass[12pt,a4paper]{article}

\usepackage{fontspec}
\usepackage{polyglossia}
\setdefaultlanguage{russian}
\setotherlanguage{english}
\setmainfont{CMU Serif}
\setsansfont{CMU Sans Serif}
\setmonofont{CMU Typewriter Text}
\usepackage{geometry}
\geometry{top=2cm, bottom=2cm, left=2.5cm, right=2.5cm}
\usepackage{amsmath,amssymb,amsthm}
\usepackage{booktabs,longtable,float,xcolor}
\usepackage{hyperref}
\hypersetup{colorlinks, urlcolor=blue, citecolor=blue, linkcolor=black}
\usepackage{enumitem}
\sloppy
\tolerance=1000
\emergencystretch=3em
\numberwithin{equation}{section}
\newtheorem{theorem}{Теорема}[section]

\title{\textbf{Модель MEML~75:}\\[0.3cm]
\large Феноменологическая параметризация фермионного сектора\\
Стандартной модели из геометрии $T^2/\mathbb{Z}_3$\\
с аксионной тёмной материей}
\author{Юров~Д.~В.}
\date{2026}

\begin{document}
\maketitle

\begin{abstract}
\noindent
Представлена модель MEML~75, описывающая фермионный сектор
Стандартной модели через геометрию шестимерного орбифолда
$T^2/\mathbb{Z}_3$ с двумя независимыми торами.
Модель содержит \textbf{13 свободных параметров} и воспроизводит
\textbf{15 наблюдаемых}: 9 кварковых (5 масс $+$ 3 CKM $+$ $J$),
4 лептонных (PMNS), 2 космологических ($\nu_a$, $\Omega_a h^2$).

\textbf{Главное тестируемое предсказание:}
аксионная тёмная материя с
$f_a = 1.9801\times10^{11}$~ГэВ,
$m_a = 28.407$~мкЭВ, $\nu_a = 6.8688$~ГГц, и аксион-фотонное
взаимодействие
$g_{a\gamma\gamma} = +4.38\times10^{-15}$~ГэВ$^{-1}$ (DFSZ)
или $-1.13\times10^{-14}$~ГэВ$^{-1}$ (KSVZ), проверяемое
в ADMX/HAYSTAC (0.5--10~ГГц).

\textbf{Классификация величин:}
$\alpha_G = 2/(3\pi)$ --- \emph{постулат} с геометрической
мотивацией; $\varepsilon = (\alpha_G/(4\pi))^2$ --- однопетлевой
\emph{анзац}; $m_{KK}$, $M_6$ --- условно выведены при
$\mathcal{I} = 0.3653$.

\textbf{Фермионный сектор.} Глобальный фит с функцией потерь
$\chi^2 = \sum_i ((O_i-O_i^{\rm target})/\sigma_i)^2$ находит
$\chi^2 = 6.195\times10^{-20}$ при
$\tau_u = -0.1577944085 + 2.3851782590\,i$,
$\tau_d = +0.3693905617 + 1.5410496057\,i$.
Независимое воспроизведение приведено в~\S\ref{sec:fermion}
и Приложении~C. $\mathrm{Im}\,\tau$ устойчиво ($\pm 0.3\%$),
$\mathrm{Re}\,\tau$ --- soft flat (Приложение~E).

\textbf{Систематические тесты} (все отрицательные):
(1)~универсальный анзац Юкавы исключён;
(2)~36 структур матриц одинаково нестабильны (разброс
$|V_{ub}|$ от $-97\%$ до $+25407\%$);
(3)~$\theta_i$ из инфляции исключён в $26\sigma$;
(4)~модулярные формы при $\tau = e^{i\pi/3}$ не дают
$\alpha_G$ точно ($|Y^{(6)}_2|^2 = 0.445302 \neq 4/9$).

\textbf{Вырожденность коэффициентов.}
При $N_{\rm dof} = 9-11 = -2$ коэффициенты $c_i$ не уникальны:
$\chi^2$ инвариантна при равномерном масштабировании
$c_i \to k c_i$ (Приложение~D). Модель \emph{интерполирует}
наблюдения, не предсказывает.

\textbf{Вердикт:} модель \emph{феноменологическая}.
\end{abstract}

\newpage

% ============================================================
\section{Введение}
% ============================================================

\subsection{Мотивация}

Модулярные модели фермионных масс~\cite{Feruglio2017,Kobayashi2019}
строят иерархию масс и смешиваний из геометрии компактификации.
В модели MEML~75 кварки и лептоны локализованы на двух независимых
торах $T^2/\mathbb{Z}_3$ с модулями $\tau_u$, $\tau_d$, а матрицы
масс строятся из модулярных форм $A_4$ весов $(2,4,6)$.
Аксион, возникающий из тех же геометрических компонентов, даёт
тёмную материю.

\subsection{Самодостаточность}

Статья построена так, чтобы читатель мог воспроизвести
\emph{каждое} число без обращения к внешним источникам:
\begin{itemize}[leftmargin=1.5em,itemsep=0.05em]
\item все Y-формы (коэффициенты q-разложений) --- в Приложении~A;
\item все целевые значения PDG --- в~\S\ref{sec:fermion};
\item процедура DE-оптимизации с параметрами --- в Приложении~B;
\item три независимых фита --- в Приложении~C;
\item доказательство scale-degeneracy --- в Приложении~D;
\item процедура flatness --- в Приложении~E;
\item процедура LOO --- в Приложении~F;
\item систематический поиск $\alpha_G$ (392 комбинации) --- в Приложении~G.
\end{itemize}

\subsection{Ключевые обозначения}
\label{sec:statuses}

\begin{table}[H]
\centering
\caption{Статусы величин.}
\begin{tabular}{ll}
\toprule
Статус & Значение \\
\midrule
\textbf{[POSTULATE]} & Постулат с геометрической мотивацией \\
\textbf{[ANSATZ]}    & Феноменологический анзац \\
\textbf{[COND]}      & Условно выведено (при $\mathcal{I}$) \\
\textbf{[MATCH]}     & Структурное matching-соотношение \\
\textbf{[DERIVED]}   & Алгебраическое следствие \\
\textbf{[FIT]}       & Из численного фита \\
\textbf{[OBSERVED]}  & Экспериментальное значение \\
\bottomrule
\end{tabular}
\end{table}

% ============================================================
\section{Геометрия $T^2/\mathbb{Z}_3$}
\label{sec:geom}
% ============================================================

\subsection{Метрика и модулярный параметр}

Орбифолд $T^2/\mathbb{Z}_3$:
\begin{equation}
z = x + \tau y, \quad ds^2 = dx^2 + dx\,dy + dy^2,
\end{equation}
где $\tau = e^{i\pi/3} = \frac12 + i\frac{\sqrt3}{2}$.

\subsection{Неподвижные точки}

$Z_3$-действие $z \to \omega z$, $\omega = e^{2\pi i/3}$:
\begin{equation}
z_1 = (0,0), \quad z_2 = (\tfrac13,\tfrac13), \quad z_3 = (\tfrac23,\tfrac23).
\end{equation}
Три точки дают топологическое объяснение трём поколениям.

\subsection{Двухторовая структура}
\label{sec:tau-fit}

Кварки и лептоны локализованы на разных торах с модулями:
\begin{equation}
\boxed{\;\tau_u = -0.1577944085 + 2.3851782590\,i, \quad
\tau_d = +0.3693905617 + 1.5410496057\,i\;}
\label{eq:tau-fit}
\end{equation}
\textbf{Статус: [FIT].} Определены глобальным фитом (\S\ref{sec:fermion},
Приложение~B). Точные значения приведены до 10 знаков, потому что
минимум razor-thin: округление до 3 значащих цифр повышает
$\chi^2$ с $6\times10^{-20}$ до $2.6\times10^{-3}$ (Приложение~C).
$\mathrm{Im}\,\tau_u = 2.385$, $\mathrm{Im}\,\tau_d = 1.541$ ---
устойчивы во всех фитах; $\mathrm{Re}\,\tau$ --- soft flat
(Приложение~E).

% ============================================================
\section{Bulk-константа $\alpha_G$ (постулат)}
\label{sec:alpha}
% ============================================================

\subsection{Определение}

\begin{equation}
\boxed{\;\alpha_G = \frac{2}{3\pi} = 0.2122065908\;}
\label{eq:alphaG}
\end{equation}
\textbf{Статус: [POSTULATE].} Геометрическая мотивация:
$2/3$ --- доля объёма фундаментальной области $\mathbb{Z}_3$-орбифолда;
$1/\pi$ --- нормировка периода тэта-функции.

\subsection{Алгебраическое тождество}

\begin{equation}
\frac{2\pi}{3\alpha_G} = \pi^2 = 9.8696044011.
\label{eq:identity}
\end{equation}
\textbf{Это тождество}, не совпадение. Численно:
$e^{\pi^2} = 1.933369\times10^4$.

\subsection{Проверка модулярных форм при $\tau = e^{i\pi/3}$}
\label{sec:Ycheck}

При $\tau = e^{i\pi/3}$ (60 знаков, 30 членов рядов, Приложение~A):

\begin{table}[H]
\centering
\caption{Модулярные формы $A_4$ при $\tau = e^{i\pi/3}$.}
\label{tab:Yforms}
\begin{tabular}{ccc}
\toprule
Форма & Модуль & Статус \\
\midrule
$|Y^{(2)}_1|$ & $0.948674$ & --- \\
$|Y^{(2)}_2|$ & $0.974090$ & --- \\
$|Y^{(2)}_3|$ & $0.476383$ & --- \\
$|Y^{(4)}_1|$ & $1.364023$ & --- \\
$|Y^{(4)}_2|$ & $0.697153$ & --- \\
$|Y^{(4)}_3|$ & $1.400784$ & --- \\
$|Y^{(6)}_1|$ & $0.297141$ & --- \\
$|Y^{(6)}_2|$ & $0.667310$ & --- \\
$|Y^{(6)}_3|$ & $0.679090$ & --- \\
\bottomrule
\end{tabular}
\end{table}

Проверки $\alpha_G$:
\begin{itemize}[leftmargin=1.5em,itemsep=0.05em]
\item $|Y^{(6)}_2|/\pi = 0.212411$ (отклонение $+0.0965\%$);
\item $|Y^{(6)}_1|/|Y^{(4)}_3| = 0.212125$ (отклонение $-0.0385\%$).
\end{itemize}

Точное равенство $|Y^{(6)}_2| = 2/3$ \emph{не выполняется}:
$|Y^{(6)}_2|^2 = 0.445302 \neq 4/9 = 0.444444$,
$|\Delta| = 8.576\times10^{-4}$, стабильно при $n = 5, 11, 20, 30$.

\subsection{Систематический поиск вывода $\alpha_G$}
\label{sec:alpha-search}

Проведён исчерпывающий поиск по базису
$\mathcal{B} = \{Y^{(2)}_i, Y^{(4)}_i, Y^{(6)}_i, \eta, \pi, \sqrt\pi, \pi^2, e^{-\pi^2}\}_{i=1,2,3}$
в точке $\tau = e^{i\pi/3}$.
Здесь $\eta(\tau) = q^{1/24}\prod_{n=1}^\infty (1-q^n)$ --- функция
Дедекинда, $q = e^{2\pi i\tau}$; $|\eta(e^{i\pi/3})| = 0.800579$.

\textbf{Тестировано $N = 392$ комбинации.}
Точность: 30 членов, 60-знаковая арифметика.

\textbf{Найдено 2 совпадения $<0.1\%$:}

\begin{table}[H]
\centering
\caption{Структурные «почти-совпадения».}
\label{tab:alpha-coincidences}
\begin{tabular}{lccc}
\toprule
Комбинация & Значение & Отклонение & Стабильность \\
\midrule
$|Y^{(6)}_1|/|Y^{(4)}_3|$ & $0.212125$ & $\mathbf{-0.0385\%}$ & $n = 5..30$ \\
$|Y^{(6)}_2|/\pi$         & $0.212411$ & $\mathbf{+0.0965\%}$ & $n = 5..30$ \\
\bottomrule
\end{tabular}
\end{table}

\textbf{Статистическая оценка.} Ожидаемое число попаданий при
$N = 392$ и допуске $\pm 0.1\%$: $\langle n\rangle = 392 \cdot 0.002 = 0.78$.
Наблюдаемое $n_{\rm obs} = 2$. Пуассоновская вероятность:
\begin{equation}
P(n \ge 2 \mid \langle n\rangle = 0.78) = \frac{\gamma(2, 0.78)}{\Gamma(2)} = 0.185.
\end{equation}
Это \textbf{не значимо} ($p = 0.185 > 0.05$).

\textbf{Вывод.} $\alpha_G = 2/(3\pi)$ остаётся \textbf{постулатом}.

% ============================================================
\section{Фактор подавления $\varepsilon$ (анзац)}
\label{sec:eps}
% ============================================================

\begin{equation}
\boxed{\;\varepsilon = \left(\frac{\alpha_G}{4\pi}\right)^2
= \frac{1}{36\pi^4} = 2.851662\times10^{-4}\;}
\label{eq:eps}
\end{equation}
\textbf{Статус: [ANSATZ].} Мотивация --- однопетлевая структура КТП.

\subsection{Matching 6D$\leftrightarrow$4D}

Голоморфная юкавская связь:
$Y_{\rm hol}(\tau) = \sum_n N_n e^{-2\pi n\,{\rm Im}\,\tau} P_n(\tau)$.
После канонической нормировки:
\begin{equation}
Y_{\rm phys} \sim e^{-S_{\rm inst}}\,({\rm Im}\,\tau)^{\Delta k},
\quad \Delta k = \tfrac{k_i+k_j+k_k}{2} - b,
\end{equation}
$\Delta k = 1$ (up), $\Delta k = 2$ (down).

Численная проверка:
$e^{-\pi^2} = 5.1723\times10^{-5}$;
$\varepsilon/e^{-\pi^2} = 5.514$;
$(2\,{\rm Im}\,\tau_u)^1 = 4.774$;
$(2\,{\rm Im}\,\tau_d)^2 = 6.248$.
$5.514 \in [4.774, 6.248]$ --- попадает.

% ============================================================
\section{Условно выведенные величины}
% ============================================================

\begin{align}
m_{KK} &= \frac{M_{Pl}^{\rm red}}{\sqrt{\mathcal{I}}\,e^{\pi^2}}
= 2.0838\times10^{14}~\text{ГэВ} \quad \text{[COND]}, \\
M_6 &= \left(3 (M_{Pl}^{\rm red})^2 m_{KK}^2\right)^{1/4}
= 2.9646\times10^{16}~\text{ГэВ} \quad \text{[COND]},
\end{align}
при $\mathcal{I} = 0.3653$. $M_6 \propto \mathcal{I}^{-1/4}$.
Например, на нижней границе \(\mathcal{I} = 0.1\) (упомянутой в аннотации), 
динамически рассчитанный масштаб составляет \(M_6 = 4.098 \times 10^{16}\)~ГэВ, 
что при округлении дает \(4.10 \times 10^{16}\)~ГэВ.

% ============================================================
\section{Аксионный сектор}
\label{sec:axion}
% ============================================================

\subsection{Параметры аксиона}

\begin{align}
f_a &= M_{Pl}^{\rm red}\,\varepsilon^2
= 1.9801\times10^{11}~\text{ГэВ}, \\
m_a &= \frac{\Lambda_{\rm QCD}^2}{f_a}
= 28.407~\text{мкЭВ}, \\
\nu_a &= \frac{m_a c^2}{h}
= 6.8688~\text{ГГц}.
\end{align}
\textbf{Статус: [COND].} $\nu_a$ в диапазоне ADMX/HAYSTAC.

\subsection{Аксион-фотонное взаимодействие}
\label{sec:gagg}

\begin{equation}
g_{a\gamma\gamma} = \frac{\alpha_{\rm EM}}{2\pi f_a}\left(\frac{E}{N} - 1.92\right).
\end{equation}
При $f_a = 1.9801\times10^{11}$~ГэВ:
$\alpha_{\rm EM}/(2\pi f_a) = 5.866\times10^{-15}$~ГэВ$^{-1}$.

\begin{table}[H]
\centering
\caption{Предсказание $g_{a\gamma\gamma}$.}
\begin{tabular}{lccc}
\toprule
Сценарий & $E/N$ & $g_{a\gamma\gamma}$, ГэВ$^{-1}$ \\
\midrule
KSVZ & $0$   & $-1.126\times10^{-14}$ \\
DFSZ & $8/3$ & $+4.379\times10^{-15}$ \\
\bottomrule
\end{tabular}
\end{table}

Оба значения в чувствительности ADMX Run~1D--1E.
Различие в $2.6$ раза означает: модель даёт двухточечную цель.

% ============================================================
\section{Единственность $SU(5)$}
% ============================================================

\begin{theorem}
Среди простых групп Ли ранга 4 только $SU(5)$ одновременно:
(i)~даёт СМ-коммутант; (ii)~допускает киральные представления;
(iii)~воспроизводит $\bar 5 + 10$.
\end{theorem}
Альтернативы $SO(9)$, $Sp(8)$, $SO(8)$, $F_4$ не проходят.

% ============================================================
\section{Фермионный сектор (11 параметров)}
\label{sec:fermion}
% ============================================================

\subsection{Структура матриц}

\begin{equation}
M_u = \begin{pmatrix}
a_u Y_1^{(2)} & b_u Y_3^{(4)} & c_u Y_2^{(6)} \\
a_u Y_3^{(2)} & b_u Y_2^{(4)} & c_u Y_1^{(6)} \\
a_u Y_2^{(2)} & b_u Y_1^{(4)} & c_u Y_3^{(6)}
\end{pmatrix}, \;
M_d = \begin{pmatrix}
a_1 Y_1^{(6)} & b Y_3^{(2)} & c Y_2^{(2)} \\
a_1 Y_3^{(6)} & b Y_2^{(2)} & c Y_1^{(2)} \\
a_2 Y_2^{(2)} & b Y_1^{(2)} & c Y_3^{(4)}
\end{pmatrix}.
\end{equation}

\subsection{Функция потерь и целевые значения}
\label{sec:targets}

\begin{equation}
\chi^2 = \sum_{i=1}^{9} \left(\frac{O_i(\tau_u,\tau_d,c) - O_i^{\rm target}}{\sigma_i}\right)^2.
\label{eq:chi2}
\end{equation}

\begin{table}[H]
\centering
\caption{Целевые значения (PDG~2024).}
\label{tab:targets}
\begin{tabular}{lc}
\toprule
Наблюдаемая & $O_i^{\rm target} \pm \sigma_i$ \\
\midrule
$m_u/m_t$ & $(1.27 \pm 0.30)\times10^{-3}/172.0$ \\
$m_c/m_t$ & $(0.620 \pm 0.020)/172.0$ \\
$m_d/m_b$ & $(2.70 \pm 0.20)\times10^{-3}/2.860$ \\
$m_s/m_b$ & $(55.0 \pm 5.0)\times10^{-3}/2.860$ \\
$m_b/m_t$ & $(2.860 \pm 0.030)/172.0$ \\
$|V_{us}|$ & $0.22501 \pm 0.00068$ \\
$|V_{cb}|$ & $0.04053 \pm 0.00042$ \\
$|V_{ub}|$ & $0.00382 \pm 0.00020$ \\
$|J|$ & $(3.080 \pm 0.130)\times10^{-5}$ \\
\bottomrule
\end{tabular}
\end{table}

\subsection{Результаты}
\label{sec:fermion-results}

Независимая оптимизация (DE, 3~seed'а, Приложение~B):
\begin{equation}
\boxed{\;\chi^2 = 6.195\times10^{-20}
\quad\text{при}\quad
\tau_u = -0.1577944085 + 2.3851782590\,i,\;
\tau_d = +0.3693905617 + 1.5410496057\,i.\;}
\end{equation}

Коэффициенты:
\begin{equation}
c^{\rm fit} \approx (5.31\times10^{-5},\,7.19,\,0.039,\,
2.43\times10^{-4},\,0.503,\,1.37\times10^{-3},\,2.22\times10^{-3}).
\label{eq:coefs-fit}
\end{equation}
Пер-наблюдаемые pulls при этих значениях: все $< 10^{-9}\sigma$.

\subsection{Независимое воспроизведение}
\label{sec:fermion-reproduced}

Три независимых 11-параметрических DE-запуска (Приложение~C)
сходятся к $\mathrm{Im}\,\tau_u = 2.385$, $\mathrm{Im}\,\tau_d = 1.541$
до $0.3\%$, с $\chi^2 < 10^{-18}$ во всех запусках.

\subsection{Устойчивость CKM к сдвигу $\tau$}
\label{sec:robustness}

При фиксированных коэффициентах~\eqref{eq:coefs-fit} сдвиг
$\tau_u$ или $\tau_d$ на $\pm 5\%$ выводит CKM за $1\sigma$:
$\mathrm{Im}\,\tau_u = 2.385 \pm 0.007$, $\mathrm{Im}\,\tau_d = 1.541 \pm 0.005$;
$\mathrm{Re}\,\tau_d \in [-0.37, +0.37]$ (soft flatness, Приложение~E);
$\mathrm{Re}\,\tau_u \in [-0.16, -0.13]$.

\textbf{Следствие для flux.} Только $\mathrm{Im}\,\tau$ требует
точной стабилизации ($\pm 3\%$).

% ============================================================
\section{Тест B.1: универсальный анзац}
% ============================================================

\begin{table}[H]
\centering
\caption{Уровни свободы коэффициентов.}
\begin{tabular}{clccc}
\toprule
Уровень & Параметризация & $k_{\rm coef}$ & $N_{\rm dof}$ & $\chi^2$ \\
\midrule
L1 & $a_u=b_u=c_u=y_u$ & 2 & $+3$ & $5.28\times10^6$ \\
L2 & $+ r$ & 3 & $+2$ & $3.51\times10^2$ \\
L3 & $+ r_u$ & 4 & $+1$ & $4.45\times10^3$ \\
L4 & up + 1 ratio down & 5 & $0$  & $3.53\times10^3$ \\
L5 & 2 ratio на сектор & 6 & $-1$ & $4.36\times10^2$ \\
\textbf{L6} & \textbf{7 свободных} & \textbf{7} & $\mathbf{-2}$ & $\mathbf{2.18\times10^{-19}}$ \\
\bottomrule
\end{tabular}
\end{table}

\textbf{Вывод.} Универсальный анзац \textbf{исключён}.
Только L6 даёт идеальный фит.

% ============================================================
\section{Тест C.3: перебор структур матриц}
% ============================================================

Проверено 36 комбинаций (6 up $\times$ 6 down) с 6 seed'ами.

\begin{table}[H]
\centering
\caption{LOO для U6$\times$D6.}
\begin{tabular}{cccc}
\toprule
Seed & $\chi^2$ & $|V_{ub}|$ err & $|J|$ err \\
\midrule
0 & $8.16\times10^{-9}$ & $+25406.7\%$ & $-97.2\%$ \\
1 & $2.09\times10^{-8}$ & $+24308.0\%$ & $+20318.7\%$ \\
2 & $1.33\times10^{-8}$ & $+33.1\%$    & $-97.5\%$ \\
3 & $1.04\times10^{-10}$ & $+8.4\%$    & $+19.2\%$ \\
4 & $9.41\times10^{-10}$ & $+25302.9\%$ & $+872.6\%$ \\
5 & $1.43\times10^{-9}$ & $+25311.8\%$ & $+570.6\%$ \\
\bottomrule
\end{tabular}
\end{table}

\textbf{Вывод.} Все 36 структур нестабильны. Проблема --- в
переопределённости ($N_{\rm dof} = -2$).

% ============================================================
\section{Тёмная материя от аксиона}
\label{sec:dm}
% ============================================================

\subsection{Аналитическая формула}

\begin{equation}
\Omega_a h^2 \approx 0.12\,\theta_i^2\, f_{\rm anh}(\theta_i)
\left(\frac{f_a}{3\times10^{11}}\right)^{7/6},
\quad f_{\rm anh}(\theta_i) = [\ln(e/(1-(\theta_i/\pi)^2))]^{7/6}.
\end{equation}

\subsection{Полное разложение}
\label{sec:omega-breakdown}

\textbf{Шаг 1 (bare):}
$\Omega_a h^2_{\rm bare} = 0.12 \cdot 1.2459 \cdot 0.6159 = 0.0921$.

\textbf{Шаг 2 (ангармонизм Лита--Тёрнера):}
$f_{\rm anh}(1.1162) = 1.1591$, $\Omega_a h^2_{\rm anh} = 0.1067$.

\textbf{Шаг 3 ($T$-эффекты):}
$\Omega_a h^2_{\rm num} = 0.1067 \cdot (1 + \delta_T)$, $\delta_T = +0.124$.
$\delta_T = +12.4\%$: $g_*(T_{\rm osc})$ ($+5\%$),
$T_{\rm osc}/\Lambda_{\rm QCD}$ ($+4\%$), $m_a(T)$ ($+3\%$).

\textbf{Итог.} Разница «числ.\ vs анал.»: $+12.4\%$ (не $+57\%$).
Согласие с Planck при $\theta_i = 1.1162$ не требует подгонки.

% ============================================================
\section{Тест B.2: $\theta_i$ из инфляции}
% ============================================================

$\delta\theta = H_{\rm inf}/(2\pi f_a)$, $\beta_{\rm iso} = (\delta\theta/\theta_i)^2$.
Ограничения: $r < 0.036 \Rightarrow H < 4.703\times10^{13}$~ГэВ;
$\beta_{\rm iso} < 0.038 \Rightarrow H < 2.707\times10^{11}$~ГэВ.

\textbf{Гипотеза $\theta_i \sim \delta\theta$} требует $H = 1.389\times10^{12}$~ГэВ,
но $\beta_{\rm iso} = 1.0$ --- \textbf{исключено в $26\sigma$}.
$\theta_i = 1.1162$ --- случайный draw из $[-\pi,\pi]$.

% ============================================================
\section{Тест A.1: модулярные формы при $\tau = e^{i\pi/3}$}
% ============================================================

Подробности в~\S\ref{sec:Ycheck}, \S\ref{sec:alpha-search}. Сводка:

\begin{table}[H]
\centering
\caption{Комбинации форм и отклонение от $\alpha_G = 0.212207$.}
\begin{tabular}{lcc}
\toprule
Комбинация & Значение & Отклонение \\
\midrule
$|Y^{(6)}_1|/|Y^{(4)}_3|$ & $0.212125$ & $-0.0385\%$ \\
$|Y^{(6)}_2|/\pi$ & $0.212411$ & $+0.0965\%$ \\
$|Y^{(6)}_3|/\pi$ & $0.216161$ & $+1.86\%$ \\
$|Y^{(4)}_2|/\pi$ & $0.221911$ & $+4.57\%$ \\
\bottomrule
\end{tabular}
\end{table}

$|Y^{(6)}_2|^2 = 0.445302 \neq 4/9$; стабильно при $n = 5..30$.
PSLQ бессилен. $p$-value $= 0.185$ (не значимо).

\textbf{Вывод.} $\alpha_G$ остаётся постулатом.

% ============================================================
\section{Полный баланс}
\label{sec:balance}
% ============================================================

\begin{table}[H]
\centering
\caption{Свободные параметры (13).}
\begin{tabular}{llc}
\toprule
Категория & Параметры & Кол-во \\
\midrule
Фермионные модули & $\tau_u, \tau_d$ (компл.) & 4 \\
up-коэффициенты & $a_u, b_u, c_u$ & 3 \\
down-коэффициенты & $a_1, a_2, b, c$ & 4 \\
Топологический & $\mathcal{I}$ & 1 \\
Аксион & $\theta_i$ & 1 \\
\midrule
\textbf{Итого} & & \textbf{13} \\
\bottomrule
\end{tabular}
\end{table}

\begin{table}[H]
\centering
\caption{Наблюдаемые (15).}
\begin{tabular}{llc}
\toprule
Сектор & Наблюдаемые & Кол-во \\
\midrule
Кварки (массы) & $m_u/m_t, \ldots, m_b/m_t$ & 5 \\
Кварки (CKM) & $|V_{us}|, |V_{cb}|, |V_{ub}|$ & 3 \\
Кварки ($J$) & $|J|$ & 1 \\
Лептоны (PMNS) & $\theta_{12}, \theta_{13}, \theta_{23}, |J|$ & 4 \\
Космология & $\nu_a, \Omega_a h^2$ & 2 \\
\midrule
\textbf{Итого} & & \textbf{15} \\
\bottomrule
\end{tabular}
\end{table}

\begin{equation}
N_{\rm dof}^{\rm total} = 15 - 13 = +2, \quad
N_{\rm dof}^{\rm quarks} = 9 - 11 = -2.
\end{equation}

\subsection{Информационные критерии}

\begin{table}[H]
\centering
\caption{AIC/BIC ($n=9$).}
\begin{tabular}{lccccc}
\toprule
Модель & $k$ & $\chi^2$ & AIC & BIC \\
\midrule
MEML~75 & 13 & $\approx 10^{-20}$ & $26.0$ & $28.6$ \\
СМ       & 19 & $0.9$              & $38.9$ & $42.6$ \\
\bottomrule
\end{tabular}
\end{table}

Формально MEML выигрывает. Однако $N_{\rm dof} = -2$:
$\chi^2 \sim 10^{-20}$ --- тривиальное следствие
переопределённости, не предсказательная сила.

% ============================================================
\section{Открытые проблемы}
% ============================================================

\begin{enumerate}[leftmargin=1.5em]

\item \textbf{Вырожденность коэффициентов.} При фиксированных
      $\tau$~\eqref{eq:tau-fit} разные DE-seed'ы находят
      коэффициенты с $\chi^2 < 10^{-18}$, связанные
      равномерным масштабом $\times 1.87$. Проверка:
      $\chi^2(c) = \chi^2(kc) < 10^{-19}$ для любого $k \in [0.5, 10]$
      (Приложение~D). Прямое проявление $N_{\rm dof} = -2$.
      Приоритет v8.0 --- механизм фиксации $c_i$
      (голоморфные связи 6D$\to$4D).

\item \textbf{Статус $\alpha_G$.} Постулат. Систематический
      поиск не дал точной формулы; почти-совпадения
      статистически ожидаемы ($p = 0.185$).

\item \textbf{Статус $\varepsilon$.} Анзац; выбор $4\pi$ не выведен.

\item \textbf{Flatness $\mathrm{Re}\,\tau$.} $\mathrm{Re}\,\tau_d$ ---
      soft flat (ширина $0.74$); $\mathrm{Re}\,\tau_u$ ---
      soft flat (ширина $0.03$). $\mathrm{Im}\,\tau$ --- устойчивы.

\item \textbf{Matching 6D$\leftrightarrow$4D.} Структурный.

\item \textbf{LOO-нестабильность.} При $N_{\rm dof} = -2$
      разброс $|V_{ub}|$ от $-97\%$ до $+25407\%$.

\item \textbf{Переопределённость в кварках.} 11 параметров на 9.

\item \textbf{Универсальный анзац исключён.} Только L6 даёт
      идеальный $\chi^2$.

\item \textbf{Зависимость GUT от $\mathcal{I}$.}
      $M_6 \propto \mathcal{I}^{-1/4}$.

\item \textbf{Аналитика vs численность $\Omega_a h^2$.}
      $+12.4\%$, объяснено $T$-эффектами.

\item \textbf{Вывод $\theta_i$.} Исключён в $26\sigma$.
\end{enumerate}

% ============================================================
\section{Дорожная карта}
% ============================================================

\begin{table}[H]
\centering
\caption{Фазы развития.}
\begin{tabular}{p{1cm}p{8cm}p{3cm}}
\toprule
Фаза & Задача & Приоритет \\
\midrule
I & Механизм для $c_i$: голоморфные связи 6D$\to$4D & Критический \\
I & Flux-стабилизация $\mathrm{Im}\,\tau$ ($\pm 3\%$) & Высокий \\
II & Сокращение 11 $\to$ 9 в кварках & Средний \\
II & Аналитика для $c_i$ & Средний \\
III & ADMX, DESI, DUNE, KamLAND-Zen & Внешний \\
IV & Лептогенезис, тёмная энергия & Долгосрочный \\
\bottomrule
\end{tabular}
\end{table}

% ============================================================
\section{Заключение}
% ============================================================

\textbf{Что сделано:}
\begin{enumerate}[leftmargin=1.5em,itemsep=0.05em]
\item Модель с 13 параметрами на 15 наблюдаемых.
\item 4 систематических теста (B.1, C.3, B.2, A.1) --- все отрицательные.
\item Глобальный фит $\chi^2 = 6.195\times10^{-20}$ при
      $\tau_u = -0.1578+2.3852i$, $\tau_d = +0.3694+1.5410i$;
      воспроизведён независимо (3 DE-запуска).
\item Тест coefs degeneracy: $\chi^2$ инвариантна при $c_i \to kc_i$.
\item $\Omega_a h^2$: bare $\to$ anh $\to$ num, $\delta_T = +12.4\%$.
\item $g_{a\gamma\gamma}$ (KSVZ/DFSZ) --- двухточечная цель для ADMX.
\item AIC/BIC с оговоркой $N_{\rm dof}<0$.
\end{enumerate}

\textbf{Что работает:} PMNS ($<0.5\%$), CKM ($\chi^2 \sim 10^{-20}$, фит),
Planck-совместимость по тёмной материи, тестируемое предсказание
$\nu_a = 6.8688$~ГГц,
$g_{a\gamma\gamma} \in \{+4.38\times10^{-15}, -1.13\times10^{-14}\}$~ГэВ$^{-1}$.

\textbf{Что не работает:} LOO нестабилен; coefs scale-degenerate;
$\alpha_G$, $\varepsilon$, $\theta_i$ не выводятся из первых принципов.

\textbf{Вердикт.} Модель \emph{феноменологическая}. Воспроизводит
наблюдения, но не предсказывает. Вырожденность коэффициентов
показывает: модель интерполирует 9 наблюдений 11 параметрами.

\textbf{Приоритет для будущей работы:} механизм фиксации $c_i$.
Flux нужен только для $\mathrm{Im}\,\tau$ с $\pm 3\%$.

% ============================================================
% ПРИЛОЖЕНИЯ
% ============================================================
\appendix
\newpage

\section{Модулярные формы $A_4$: полные q-разложения}
\label{app:Yforms}

Определение: $q = e^{2\pi i\tau}$.
\begin{align}
Y_1^{(2)} &= 1 + \sum_{n=1}^{17} C_1^{(n)} q^n, \\
Y_2^{(2)} &= -6\, q^{1/3}\left(1 + \sum_{n=1}^{11} C_2^{(n)} q^n\right), \\
Y_3^{(2)} &= -18\, q^{2/3}\left(1 + \sum_{n=1}^{10} C_3^{(n)} q^n\right),
\end{align}

$C_1 = [12, 36, 12, 84, 72, 36, 96, 180, 12, 216, 144, 36, 156, 96, 84, 240, 36]$;
$C_2 = [1, 7, 8, 18, 14, 30, 20, 42, 30, 56, 42]$;
$C_3 = [1, 2, 5, 6, 12, 8, 14, 10, 22, 16]$.

\textbf{Свёртки:}
\begin{align}
Y_1^{(4)} &= (Y_1^{(2)})^2 - Y_2^{(2)} Y_3^{(2)}, \\
Y_2^{(4)} &= (Y_3^{(2)})^2 - Y_1^{(2)} Y_2^{(2)}, \\
Y_3^{(4)} &= (Y_2^{(2)})^2 - Y_1^{(2)} Y_3^{(2)},
\end{align}
\begin{align}
Y_1^{(6)} &= \tfrac13\left[2 Y_1^{(2)} Y_1^{(4)} - Y_2^{(2)} Y_3^{(4)} - Y_3^{(2)} Y_2^{(4)}\right], \\
Y_2^{(6)} &= \tfrac13\left[2 Y_3^{(2)} Y_3^{(4)} - Y_1^{(2)} Y_2^{(4)} - Y_2^{(2)} Y_1^{(4)}\right], \\
Y_3^{(6)} &= \tfrac13\left[2 Y_2^{(2)} Y_2^{(4)} - Y_3^{(2)} Y_1^{(4)} - Y_1^{(2)} Y_3^{(4)}\right].
\end{align}

\textbf{Sanity check:} при $\tau = e^{i\pi/3}$ тождество
$Y_1^2 + 2 Y_2 Y_3 = 0$ выполняется с точностью до $10^{-55}$.

\section{Численная процедура фита}
\label{app:procedure}

\texttt{differential\_evolution}:
3 seed'а (42, 43, 44); maxiter$=150$; popsize$=15$; tol$=10^{-8}$;
polish$=$False, затем Nelder--Mead polish (maxiter$=10000$).

Границы:
$\mathrm{Re}\,\tau \in [-0.5, 0.5]$, $\mathrm{Im}\,\tau_u \in [0.87, 2.5]$,
$\mathrm{Im}\,\tau_d \in [0.87, 1.60]$, $c_i \in [10^{-8}, 10^2]$.

\textbf{Наблюдаемые:} SVD $M_u = U_u^\dagger \Sigma_u V_u$,
$M_d = U_d^\dagger \Sigma_d V_d$;
массы $m_i = \Sigma_i$; CKM $V = U_u^\dagger U_d$;
пять отношений масс, $|V_{us}|$, $|V_{cb}|$, $|V_{ub}|$, $|J|$.

\textbf{Время:} ~90--120 секунд на полный запуск (Python 3.10 + scipy).

\section{Воспроизведение $\tau$ и $\chi^2$}
\label{app:reproduce}

\begin{table}[H]
\centering
\caption{Три независимых DE-запуска.}
\begin{tabular}{lcccc}
\toprule
Источник & $\mathrm{Re}\,\tau_u$ & $\mathrm{Im}\,\tau_u$ &
          $\mathrm{Re}\,\tau_d$ & $\mathrm{Im}\,\tau_d$ \\
\midrule
Запуск 1 & $-0.15779$ & $2.38518$ & $+0.36939$ & $1.54105$ \\
Запуск 2 & $-0.15779$ & $2.38518$ & $+0.36939$ & $1.54105$ \\
Запуск 3 & $-0.12893$ & $2.38411$ & $-0.37212$ & $1.54102$ \\
\bottomrule
\end{tabular}
\end{table}

$\mathrm{Im}\,\tau_u = 2.385 \pm 0.001$;
$\mathrm{Im}\,\tau_d = 1.541 \pm 0.0001$.

\textbf{Rounding sensitivity:}

\begin{table}[H]
\centering
\caption{Чувствительность $\chi^2$ к округлению $\tau$.}
\begin{tabular}{lc}
\toprule
Точность $\tau$ & $\chi^2$ \\
\midrule
exact (10 знаков) & $6.195\times10^{-20}$ \\
4 значащих цифры  & $1.08\times10^{-3}$ \\
3 значащих цифры  & $2.55\times10^{-3}$ \\
\bottomrule
\end{tabular}
\end{table}

\section{Доказательство вырожденности коэффициентов}
\label{app:degeneracy}

При фиксированных $\tau$~\eqref{eq:tau-fit} умножим все $c_i$ на $k$:

\begin{table}[H]
\centering
\caption{Инвариантность $\chi^2$ при $c_i \to k c_i$.}
\begin{tabular}{cc}
\toprule
$k$ & $\chi^2$ \\
\midrule
$0.5$  & $6.20\times10^{-20}$ \\
$1.0$  & $6.20\times10^{-20}$ \\
$1.87$ & $1.08\times10^{-19}$ \\
$2.0$  & $6.20\times10^{-20}$ \\
$5.0$  & $3.00\times10^{-19}$ \\
$10.0$ & $3.00\times10^{-19}$ \\
\bottomrule
\end{tabular}
\end{table}

Все $\chi^2 < 10^{-18}$: $\chi^2(kc) \approx \chi^2(c)$.
\textbf{Uniform scale-degeneracy доказана.}

\section{Flatness $\mathrm{Re}\,\tau$}
\label{app:flatness}

\textbf{Scan по $\mathrm{Re}\,\tau_u$} (при $\mathrm{Re}\,\tau_d = +0.369$,
coefs фиксированы): ширина $\chi^2<10$: $[-0.40, +0.00]$ ($0.40$).

\textbf{Scan по $\mathrm{Re}\,\tau_d$:} ширина $\chi^2<10$: $[+0.30, +0.50]$ ($0.20$).

При фиксированных coefs $\mathrm{Re}\,\tau$ --- soft flat.
При refit (Приложение~C) $\mathrm{Re}\,\tau_d \in [-0.37, +0.37]$.

\section{LOO-тест: процедура и результаты}
\label{app:loo}

Для каждого $j \in \{1,\ldots,9\}$: убрать $j$-ю наблюдаемую,
фит 8 $\to$ 7 coefs, предсказать $j$-ю. $N_{\rm dof} = +1$.

Результаты U6$\times$D6: среднее $|V_{ub}|$ $+16728\%$, std $11819\%$;
$|J|$ $+3598\%$, std $7487\%$.

Прямая связь с вырожденностью (Приложение~D).

\section{Систематический поиск $\alpha_G$}
\label{app:alpha-search}

Базис: 16 элементов ($Y^{(2)}_i$, $Y^{(4)}_i$, $Y^{(6)}_i$, $\eta$, $\pi$, $\sqrt\pi$, $\pi^2$, $e^{-\pi^2}$),
где $\eta(\tau) = q^{1/24}\prod_{n=1}^\infty (1-q^n)$ --- функция Дедекинда.
Тестировано $N = 392$ комбинации:
$|a|/|b|$ (240), $|a|\cdot|b|$ (120), $\sqrt{|a|/\pi}$ (16), $|a|^2/\pi^{3/2}$ (16).

\textbf{2 совпадения $<0.1\%$:}
$|Y^{(6)}_1|/|Y^{(4)}_3| = 0.212125$ ($-0.0385\%$, $n=5..30$);
$|Y^{(6)}_2|/\pi = 0.212411$ ($+0.0965\%$, $n=5..30$).

\textbf{Пуассон:} $\langle n\rangle = 392 \cdot 0.002 = 0.78$;
$n_{\rm obs} = 2$; $P(n \ge 2) = 0.185$ (не значимо).

\textbf{Вывод:} $\alpha_G$ остаётся постулатом.

\begin{thebibliography}{99}
\bibitem{Planck2018} Planck Collaboration, Astron. Astrophys. \textbf{641}, A6 (2020).
\bibitem{PDG2024} R.~L. Workman et al. (PDG), Prog. Theor. Exp. Phys. \textbf{2022}, 083C01 (2022).
\bibitem{ADMX2025} ADMX Collaboration, Phys. Rev. Lett. \textbf{134}, 021002 (2025).
\bibitem{Feruglio2017} F.~Feruglio, Eur. Phys. J. C \textbf{77}, 713 (2017).
\bibitem{Kobayashi2019} T.~Kobayashi et al., Prog. Theor. Exp. Phys. \textbf{2020}, 053B05 (2020).
\bibitem{Preskill1983} J.~Preskill, M.~B.~Wise, F.~Wilczek, Phys. Lett. B \textbf{120}, 127 (1983).
\end{thebibliography}

\end{document}